\documentclass[UFT8]{ctexart}%\setcounter{page}{1}
\setCJKmainfont[BoldFont=STHeiti,ItalicFont=STKaiti]{STSong}%\fangsong \songti  \heiti \kaishu \wuhao \small \large \LARGE
\usepackage{graphicx} % 导入插图功能
\usepackage{float}
\usepackage{threeparttable}
\usepackage{caption}
\usepackage{subfigure}
\usepackage{amsmath}
\usepackage{diagbox} % \diagbox[斜线方形]{分区1内容}{分区2内容}{分区3内容}...
\usepackage{ulem}
\usepackage{ctex}
\usepackage{booktabs}% \toprule \midrule \bottomrule
\usepackage{multirow}% \multicolumn{}{}{}
\usepackage{ulem}%\xout 斜删除线 \sout 删除线 \uline下划线 \uuline 双下划线 \uwave 波浪线 
 \usepackage{geometry} \geometry{a4paper}
\usepackage{graphicx}

\newcommand\mydot[1]{\scalebox{#1}{.}}
%\renewcommand\cftdot{\mydot{0.8}} % change the size of dots
%\renewcommand\cftdotsep{3} % change the space between dots


\usepackage{pdfpages}%\includepdf[pages=-]{插入的文件.pdf} 
%\usepackage[final]{pdfpages} %\includepdf{figure.pdf} 

\usepackage{color}
\def\red{\textcolor{red}}
\def\blue{\textcolor{blue}}
\definecolor{green}{RGB}{0,153,51}
\def\green{\textcolor{green}}

\def\bi{\begin{itemize}}
\def\ei{\end{itemize}}
\def\be{ \begin{enumerate}}
\def\ee{ \end{enumerate}}

\def\nh{\noindent\hangindent=0.45truecm\hangafter=1}

\usepackage{hyperref}% \url \href
\hypersetup{
    colorlinks=true,
    linkcolor=black,
    filecolor=blue,
    urlcolor=blue,
    pdftitle={Overleaf Example},
%    pdfpagemode=FullScreen,
    }
\urlstyle{same}
\def\link{\hyperlink}
\def\target{\hypertarget}

% 控制目录的宏包
\usepackage{titletoc}
% 按清华标准, 缩小目录中各级标题之间的缩进
\dottedcontents{section}[2cm]{\bfseries\vspace{0.4em}}{1.5em}{5pt}
\dottedcontents{subsection}[3.00cm]{}{2.7em}{5pt}
\dottedcontents{subsubsection}[4.30cm]{}{3.4em}{5pt}

%\headsep 0.5 true cm \topmargin 0pt \oddsidemargin 0pt
%\evensidemargin -10pt \textheight 297mm \textwidth 210mm

\begin{document}
%\pagestyle{empty}%全部页码为空
\thispagestyle{empty} %该页页码为空
\vspace*{3cm}
\begin{center}
{{\LARGE\heiti 研究生阶段汇报材料}\\[0.6cm]
{\normalsize 唐\ 洁}\\[0.1cm]
{\small(广西师范大学数学与统计学院, 广西桂林, 541004)}}
\end{center}

\clearpage%新的一页
%\tableofcontents%输出目录
\thispagestyle{empty} % 当前页不显示页码

\clearpage%新的一页
\setcounter{page}{1}%设置首页

\includepdf[pages=1-10]{2021-11-12.pdf}
\includepdf[pages=1-8]{2022-11-02.pdf}
\includepdf[pages=1-7]{2022-11-08.pdf}
\includepdf[pages=1-12]{2022-11-15.pdf}
\end{document}





